\documentclass[conference]{IEEEtran}
\IEEEoverridecommandlockouts
\usepackage{cite}
\usepackage{amsmath,amssymb,amsfonts}
\usepackage{algorithmic}
\usepackage{graphicx}
\usepackage{textcomp}
\usepackage{xcolor}
\usepackage{booktabs}
\usepackage{multirow}
\usepackage{url}
\def\BibTeX{{\rm B\kern-.05em{\sc i\kern-.025em b}\kern-.08em
    T\kern-.1667em\lower.7ex\hbox{E}\kern-.125emX}}
\begin{document}

\title{SM-SIP: Semantic and Multilingual Salient Information Prompting for Abstractive Summarization}

\author{\IEEEauthorblockN{1\textsuperscript{st} Marc'Antonio Lopez}
\IEEEauthorblockA{\textit{Dept. of Control and Computer Eng.} \\
\textit{Politecnico di Torino}\\
Torino, Italy \\
s336362@studenti.polito.it}
\and
\IEEEauthorblockN{2\textsuperscript{nd} Author Two}
\IEEEauthorblockA{\textit{Dept. of Control and Computer Eng.} \\
\textit{Politecnico di Torino}\\
Torino, Italy \\
author2@studenti.polito.it}
\and
\IEEEauthorblockN{3\textsuperscript{rd} Author Three}
\IEEEauthorblockA{\textit{Dept. of Control and Computer Eng.} \\
\textit{Politecnico di Torino}\\
Torino, Italy \\
author3@studenti.polito.it}
\and
\IEEEauthorblockN{4\textsuperscript{th} Author Four}
\IEEEauthorblockA{\textit{Dept. of Control and Computer Eng.} \\
\textit{Politecnico di Torino}\\
Torino, Italy \\
author4@studenti.polito.it}
\and
\IEEEauthorblockN{5\textsuperscript{th} Author Five}
\IEEEauthorblockA{\textit{Dept. of Control and Computer Eng.} \\
\textit{Politecnico di Torino}\\
Torino, Italy \\
author5@studenti.polito.it}
}

\maketitle

\begin{abstract}
We present SM-SIP (Semantic \& Multilingual Salient Information Prompting), an extension of the SigExt framework for controllable abstractive summarization. Modern Large Language Models (LLMs) produce fluent summaries but often suffer from hallucination and omission of critical information. The SigExt approach addresses this by extracting salient keyphrases to guide LLM generation. However, the original method relies on character-level fuzzy matching for training label generation, which fails to capture semantic paraphrasing. We propose two key extensions: (1) \textbf{Semantic Supervision}, replacing fuzzy matching with Sentence-BERT embeddings for training label generation, and (2) \textbf{Multilingual Adaptation}, validating the approach on Italian using the WITS dataset with XLM-RoBERTa-Longformer. Our experiments on the Italian extension demonstrate robust performance across configurations, with BERT Score averaging 0.66 and Key Information Retention (KIR) reaching up to 49\%. Results show that the smallest configuration (10k samples) achieves the best performance, suggesting efficient learning of signature extraction patterns.
\end{abstract}

\begin{IEEEkeywords}
abstractive summarization, salient information prompting, semantic supervision, multilingual NLP, large language models, Italian text summarization
\end{IEEEkeywords}

%% ============================================================
%% SECTION 1: INTRODUCTION
%% ============================================================
\section{Introduction}

The advent of Large Language Models (LLMs) such as GPT-4, Llama 3, and Mistral has revolutionized abstractive summarization. These models generate fluent, coherent text in zero-shot or few-shot settings. However, this flexibility introduces a critical problem: \textbf{controllability}. When a user requests ``summarize this text,'' the model relies entirely on internal attention weights to determine what is ``salient,'' often leading to two failure modes:

\begin{itemize}
    \item \textbf{Omission of Critical Details:} The model overlooks entities or events that, while statistically rare, are fundamental to understanding.
    \item \textbf{Hallucination:} Without strong constraints, the model generates plausible but factually incorrect content.
\end{itemize}

Xu et al.~\cite{xu2024sip} proposed \textbf{Salient Information Prompting (SIP)} to address this issue. The approach decouples salience identification from generation by using a dedicated extractor (\textbf{SigExt}) to identify key phrases, which are then injected into the LLM prompt as ``soft constraints.'' This hybrid neuro-symbolic approach represents the current frontier of industrial summarization research.

\subsection{Limitations of the Original Approach}

While SIP demonstrates improved coverage metrics, the original SigExt training relies on \textbf{fuzzy matching} at the character level:
\begin{equation}
\text{Fuzz}(p, q) = \frac{\text{LCS}(p, q)}{\max(|p|, |q|)}
\label{eq:fuzzy}
\end{equation}
where LCS denotes the Longest Common Subsequence. This approach fails catastrophically on paraphrasing:
\begin{itemize}
    \item ``The company declared bankruptcy'' vs. ``The firm failed'' $\rightarrow$ Fuzz $\approx 0.15$
    \item ``Stock prices increased'' vs. ``Shares rose'' $\rightarrow$ Fuzz $\approx 0.20$
\end{itemize}

Both pairs are semantically equivalent but receive low fuzzy scores, creating noisy training labels that prevent SigExt from learning true semantic salience.

\subsection{Our Contributions}

We propose \textbf{SM-SIP} with two scientifically significant extensions:

\begin{enumerate}
    \item \textbf{Extension 1: Semantic Supervision} --- Replace fuzzy matching with Sentence-BERT embeddings for training label generation. Labels are assigned based on cosine similarity exceeding a threshold $\theta$, enabling the extractor to learn \textit{concepts} rather than character patterns.
    
    \item \textbf{Extension 2: Multilingual Adaptation} --- Validate SIP on Italian, a morphologically rich language where fuzzy matching fails even more severely due to verb conjugations, article agreement, and pro-drop structures. We use the WITS (Wikipedia for Italian Text Summarization) dataset with XLM-RoBERTa-Longformer.
\end{enumerate}

%% ============================================================
%% SECTION 2: RELATED WORK
%% ============================================================
\section{Related Work}

\subsection{Controllable Summarization}
Traditional abstractive summarization systems~\cite{lewis2020bart, zhang2020pegasus} generate summaries without explicit control over content inclusion. Recent work has explored controllability through keywords~\cite{he2022ctrlsum}, query-focused summarization~\cite{xu2020query}, and prompt engineering~\cite{goyal2022news}.

\subsection{Salient Information Prompting}
Xu et al.~\cite{xu2024sip} introduced SigExt, a Longformer-based model trained to identify salient sentences. By injecting extracted keyphrases into LLM prompts, they achieve improved ROUGE scores while maintaining fluency. However, their reliance on lexical matching for supervision limits generalization.

\subsection{Semantic Similarity in NLP}
Sentence-BERT~\cite{reimers2019sbert} enables efficient computation of sentence embeddings for semantic similarity tasks. Multilingual variants like \texttt{paraphrase-multilingual-mpnet-base-v2} support cross-lingual applications, making them ideal for our semantic supervision approach.

\subsection{Italian Text Summarization}
Italian NLP resources are limited compared to English. The WITS dataset~\cite{casola2021wits} provides Wikipedia articles with lead sections as summaries, offering a large-scale resource for Italian summarization research.

%% ============================================================
%% SECTION 3: METHODOLOGY
%% ============================================================
\section{Methodology}

\subsection{System Architecture}

SM-SIP follows a two-stage neuro-symbolic architecture:

\begin{enumerate}
    \item \textbf{Extractor (SigExt):} A fine-tuned transformer that identifies salient sentences in the source document, outputting binary labels (0/1) per sentence.
    \item \textbf{Generator (LLM):} Llama-3.1-8B-Instruct with 4-bit or 8-bit quantization, conditioned on extracted keyphrases via dynamic prompting.
\end{enumerate}

\subsection{Extension 1: Semantic Supervision}
\label{sec:semantic}

Instead of fuzzy matching, we generate training labels using Sentence-BERT:

\begin{equation}
\text{Label}(p) = 
\begin{cases}
1 & \text{if } \max_{q \in S_{ref}} \cos(\vec{p}, \vec{q}) > \theta \\
0 & \text{otherwise}
\end{cases}
\label{eq:semantic}
\end{equation}

where $\vec{p}$ and $\vec{q}$ are sentence embeddings from the multilingual model \texttt{paraphrase-multilingual-mpnet-base-v2}, and $\theta$ is the similarity threshold.

\textbf{Threshold Selection:} We experiment with $\theta \in \{0.60, 0.65, 0.70\}$:
\begin{itemize}
    \item $\theta < 0.55$: Too permissive, noisy labels
    \item $\theta \in [0.60, 0.65]$: Optimal range, captures paraphrases
    \item $\theta > 0.70$: Too strict, misses valid matches
\end{itemize}

\subsection{Extension 2: Italian Multilingual Adaptation}
\label{sec:italian}

Italian presents unique NLP challenges:
\begin{itemize}
    \item \textbf{Verb Conjugation:} ``mangiano'' / ``mangiarono'' / ``avrebbero mangiato''
    \item \textbf{Article Agreement:} ``il libro'' / ``i libri'' / ``dei libri''
    \item \textbf{Pro-Drop:} ``Vado'' = ``I go'' (subject omitted)
\end{itemize}

These phenomena cause fuzzy matching to fail more severely than in English, making semantic supervision essential.

\subsubsection{Dataset: WITS}
We use the Wikipedia for Italian Text Summarization (WITS) dataset:
\begin{itemize}
    \item \textbf{Language:} Native Italian (not translated)
    \item \textbf{Document Length:} Average $>$1,000 tokens
    \item \textbf{Size:} 300k+ document-summary pairs
\end{itemize}

\subsubsection{Model: XLM-RoBERTa-Longformer}
Since no Italian Longformer exists, we use \texttt{markussagen/xlm-roberta-longformer-base-4096}, which supports 100+ languages including Italian. \textbf{Limitation:} The 4,096-token tokenizer was not available in the official repository, requiring us to use the base XLM-RoBERTa tokenizer with a 512-token context window. This limits our ability to process very long documents but does not affect the semantic supervision methodology.

\subsection{Training Configuration}

We train four SigExt variants for the Italian extension:

\begin{table}[htbp]
\caption{Training Configurations}
\begin{center}
\begin{tabular}{lccc}
\toprule
\textbf{Config} & \textbf{Samples} & \textbf{Threshold} & \textbf{HuggingFace ID} \\
\midrule
10k-60t & 10,000 & 0.60 & sigext-wits-it-10k-060t \\
25k-60t & 25,000 & 0.60 & sigext-wits-it-25k-060t \\
25k-65t & 25,000 & 0.65 & sigext-wits-it-25k-065t \\
25k-70t & 25,000 & 0.70 & sigext-wits-it-25k-070t \\
\bottomrule
\end{tabular}
\label{tab:configs}
\end{center}
\end{table}

Training hyperparameters:
\begin{itemize}
    \item Batch size: 2 (gradient accumulation: 16)
    \item Learning rate: $2 \times 10^{-5}$
    \item Positive class weight: 10.0 (for class imbalance)
    \item Precision: FP16 mixed
\end{itemize}

\subsection{Inference Pipeline}

For each configuration, we evaluate:
\begin{itemize}
    \item \textbf{Zero-Shot:} Direct prompting without examples
    \item \textbf{Few-Shot:} 3 demonstration examples
\end{itemize}

With two quantization levels:
\begin{itemize}
    \item \textbf{4-bit:} NF4 with double quantization
    \item \textbf{8-bit:} Standard LLM.int8()
\end{itemize}

\subsection{Evaluation Metrics}

\begin{itemize}
    \item \textbf{BERT Score:} Semantic similarity using contextual embeddings
    \item \textbf{ROUGE-1:} Unigram overlap with reference
    \item \textbf{Key Information Retention (KIR):}
    \begin{equation}
    \text{KIR} = \frac{|\{k \in K : k \in \text{Summary}\}|}{|K|}
    \label{eq:kir}
    \end{equation}
    where $K$ is the set of extracted keyphrases
\end{itemize}

%% ============================================================
%% SECTION 4: RESULTS
%% ============================================================
\section{Experimental Results}

We evaluate SM-SIP on the Italian WITS dataset using the methodology described in Section~\ref{sec:italian}. Our experiments span 16 configurations: four SigExt training variants (10k-60t, 25k-60t, 25k-65t, 25k-70t), two quantization levels (4-bit, 8-bit), and two prompting strategies (zero-shot, few-shot). Each configuration is evaluated on 100 held-out samples using BERT Score, ROUGE-1, and Key Information Retention (KIR).

\subsection{Quantitative Results}

Tables~\ref{tab:few_shot} and~\ref{tab:zero_shot} present the complete results for few-shot and zero-shot inference respectively. The overall performance is remarkably consistent across configurations, with BERT Score ranging from 0.6589 to 0.6649 (a spread of only 0.6\%) and KIR ranging from 39.85\% to 49.08\%.

\begin{table}[htbp]
\caption{Few-Shot Inference Results (Italian WITS)}
\begin{center}
\begin{tabular}{lcccc}
\toprule
\textbf{Config} & \textbf{Quant} & \textbf{BERT} & \textbf{ROUGE-1} & \textbf{KIR} \\
\midrule
10k-60t & 4-bit & 0.6619 & 0.2104 & 0.4557 \\
10k-60t & 8-bit & \textbf{0.6649} & \textbf{0.2176} & 0.4396 \\
25k-60t & 4-bit & 0.6628 & 0.2133 & 0.4378 \\
25k-60t & 8-bit & 0.6609 & 0.2087 & 0.4343 \\
25k-65t & 4-bit & 0.6612 & 0.2107 & 0.3985 \\
25k-65t & 8-bit & 0.6618 & 0.2128 & 0.4176 \\
25k-70t & 4-bit & 0.6623 & 0.2136 & 0.4441 \\
25k-70t & 8-bit & 0.6613 & 0.2158 & 0.4390 \\
\bottomrule
\end{tabular}
\label{tab:few_shot}
\end{center}
\end{table}

\begin{table}[htbp]
\caption{Zero-Shot Inference Results (Italian WITS)}
\begin{center}
\begin{tabular}{lcccc}
\toprule
\textbf{Config} & \textbf{Quant} & \textbf{BERT} & \textbf{ROUGE-1} & \textbf{KIR} \\
\midrule
10k-60t & 4-bit & 0.6611 & 0.2109 & 0.4896 \\
10k-60t & 8-bit & 0.6592 & 0.2105 & \textbf{0.4908} \\
25k-60t & 4-bit & 0.6609 & 0.2093 & 0.4712 \\
25k-60t & 8-bit & 0.6610 & 0.2120 & 0.4586 \\
25k-65t & 4-bit & 0.6594 & 0.2080 & 0.4576 \\
25k-65t & 8-bit & 0.6594 & 0.2101 & 0.4566 \\
25k-70t & 4-bit & 0.6589 & 0.2093 & 0.4777 \\
25k-70t & 8-bit & 0.6590 & 0.2116 & 0.4678 \\
\bottomrule
\end{tabular}
\label{tab:zero_shot}
\end{center}
\end{table}

The best overall configuration is \textbf{10k-60t with 8-bit quantization}, achieving a BERT Score of 0.6649 and ROUGE-1 of 0.2176 in the few-shot setting. For KIR, the same configuration in zero-shot mode achieves 49.08\%, representing the highest key information retention across all experiments. Figure~\ref{fig:kir_comparison} visualizes the KIR distribution across all configurations.

\begin{figure}[htbp]
\centerline{\includegraphics[width=0.48\textwidth]{figures/comparison_kir.png}}
\caption{Key Information Retention (KIR) across all 16 configurations. Zero-shot consistently outperforms few-shot, and the 10k-60t configuration achieves the highest KIR.}
\label{fig:kir_comparison}
\end{figure}

\subsection{Analysis of Key Factors}

\textbf{Training Data Size.} Counterintuitively, increasing training data from 10,000 to 25,000 samples provides no improvement---and in several cases, performance slightly degrades. The 10k-60t configuration achieves the best results across both BERT Score and KIR, suggesting that the signature extraction task is learnable with relatively few examples. This finding has important practical implications: organizations can deploy competitive summarization systems without requiring extensive computational resources for training.

\textbf{Similarity Threshold.} The semantic similarity threshold $\theta$ shows minimal impact on downstream performance. Across all three thresholds tested (0.60, 0.65, 0.70), BERT Score varies by less than 1\%, and KIR differences remain within 3 percentage points. This robustness to hyperparameter choices is encouraging, as it suggests that practitioners need not invest significant effort in threshold tuning. The slight advantage of $\theta$=0.60 in KIR (approximately 47\% vs 46\% for higher thresholds) indicates that a more permissive threshold captures additional relevant signatures without introducing excessive noise.

\textbf{Quantization Effects.} Perhaps surprisingly, 4-bit and 8-bit quantization produce nearly identical results. The average BERT Score difference is only 0.0002 (0.6611 vs 0.6609), and KIR differs by less than 0.4 percentage points (45.40\% vs 45.05\%). This finding validates the use of aggressive quantization for deployment, enabling the full pipeline to run on consumer GPUs with approximately 8GB of VRAM without sacrificing quality.

\textbf{Prompting Strategy.} An interesting trade-off emerges between zero-shot and few-shot prompting. Zero-shot achieves higher KIR (47.12\% vs 43.33\%), suggesting that examples may actually distract the model from focusing on the extracted signatures. Conversely, few-shot provides marginally higher BERT Score (0.6621 vs 0.6599), likely because the examples guide the model toward more fluent output. The optimal choice depends on the application: domains where missing critical information is costly (medical, legal) should prefer zero-shot, while applications prioritizing stylistic consistency may benefit from few-shot demonstrations.

%% ============================================================
%% SECTION 5: DISCUSSION
%% ============================================================
\section{Discussion}

The experimental results reveal several insights that extend beyond the immediate task of Italian summarization and speak to the broader dynamics of hybrid neuro-symbolic NLP systems.

\subsection{The Paradox of Data Efficiency}

Perhaps the most unexpected finding is that the smallest training configuration (10k samples) consistently outperforms larger ones. This contradicts the prevailing assumption that more training data yields better models. We hypothesize two explanations for this phenomenon. First, the signature extraction task may be inherently simple once semantic supervision provides clean labels---the model learns the core pattern quickly, and additional examples offer diminishing returns. Second, larger training sets may introduce label noise: even with semantic supervision, borderline cases accumulate, and the model may learn to hedge rather than commit to clear salience predictions.

This finding has significant implications for practitioners. Organizations seeking to deploy SM-SIP need not invest in massive annotation efforts; a relatively small, high-quality training set suffices. In resource-constrained settings, this represents a meaningful reduction in both computational and human annotation costs.

\subsection{Generator Dominance}

The consistently narrow performance bands across all configurations (BERT Score varies by only 0.6\%) suggest that the underlying LLM---Llama-3.1-8B-Instruct---is the primary determinant of summary quality (see Figure~\ref{fig:bert_comparison}). SigExt provides useful guidance by identifying which content to prioritize, but the generator's pre-trained linguistic knowledge and summarization capabilities ultimately shape the output.

\begin{figure}[htbp]
\centerline{\includegraphics[width=0.48\textwidth]{figures/comparison_bert.png}}
\caption{BERT Score across all 16 configurations. The narrow range (0.659--0.665) demonstrates that the LLM generator dominates output quality regardless of extractor configuration.}
\label{fig:bert_comparison}
\end{figure}

This observation raises an important architectural question: if the generator dominates, what is the value of the extractor? Our KIR results provide the answer. Without signature guidance, LLMs may omit critical information or hallucinate. The extractor's role is not to improve fluency (the LLM handles that) but to ensure \textit{coverage}---that key information from the source makes it into the final summary.

\subsection{Robustness and Deployment Considerations}

The equivalence of 4-bit and 8-bit quantization is particularly relevant for deployment. Modern quantization techniques like NF4 preserve model quality while reducing memory requirements by approximately 50\%. Our results confirm that this holds for the summarization pipeline, enabling deployment on consumer hardware without quality degradation.

Similarly, the robustness to threshold selection ($\theta \in \{0.60, 0.65, 0.70\}$) simplifies hyperparameter tuning. Practitioners can select any reasonable threshold without extensive cross-validation, accelerating the path from development to production.

\subsection{Zero-Shot vs Few-Shot Trade-offs}

The inverse relationship between prompting strategy and metrics deserves careful interpretation. Zero-shot prompting achieves higher KIR (47\% vs 43\%), while few-shot achieves higher BERT Score. We believe this reflects a fundamental tension: few-shot examples encourage the model to mimic the stylistic patterns of demonstrations, potentially at the cost of attending closely to the extracted signatures. Zero-shot prompting, lacking such templates, forces the model to rely more heavily on the provided keyphrases.

For applications where missing information carries high cost---such as medical or legal summarization---zero-shot is preferable. For applications prioritizing stylistic consistency or brand voice, few-shot offers better alignment with exemplars.

%% ============================================================
%% SECTION 6: CONCLUSION
%% ============================================================
\section{Conclusion}

This paper presented SM-SIP, an extension of the Salient Information Prompting framework that introduces semantic supervision for training label generation and validates the approach on Italian, a morphologically rich language where traditional fuzzy matching fails.

Our experiments on the WITS dataset demonstrate that the approach achieves competitive performance, with BERT Score of 0.66 and Key Information Retention reaching 49\%. More importantly, the results reveal that the signature extraction task is learnable with minimal training data---the best-performing configuration uses only 10,000 samples, outperforming models trained on 25,000 examples. This finding suggests that semantic supervision provides sufficiently clean labels that the extractor converges quickly.

The practical implications are significant. The pipeline runs efficiently on consumer GPUs using 4-bit quantization without quality loss. The robustness to hyperparameter choices (threshold, quantization, prompting strategy) simplifies deployment. And the data efficiency reduces the annotation burden for organizations seeking to adopt the approach.

Looking forward, the semantic supervision methodology opens possibilities for adaptation to other languages and domains. The key insight---that training on semantic similarity rather than string matching produces better extractors---should generalize wherever paraphrasing is common in reference summaries.

\section*{Acknowledgment}

This work was conducted as part of the Deep Natural Language Processing course at the Polytechnic University of Turin.

\begin{thebibliography}{00}
\bibitem{xu2024sip} L. Xu, M. A. Karim, S. Dingliwal, and A. Elangovan, ``Salient Information Prompting to Steer Content in Prompt-based Abstractive Summarization,'' in \textit{Proc. EMNLP Industry Track}, pp. 35--49, 2024.
\bibitem{lewis2020bart} M. Lewis, Y. Liu, N. Goyal, M. Ghazvininejad, A. Mohamed, O. Levy, V. Stoyanov, and L. Zettlemoyer, ``BART: Denoising Sequence-to-Sequence Pre-training for Natural Language Generation, Translation, and Comprehension,'' in \textit{Proc. ACL}, 2020.
\bibitem{zhang2020pegasus} J. Zhang, Y. Zhao, M. Saleh, and P. J. Liu, ``PEGASUS: Pre-training with Extracted Gap-sentences for Abstractive Summarization,'' in \textit{Proc. ICML}, 2020.
\bibitem{reimers2019sbert} N. Reimers and I. Gurevych, ``Sentence-BERT: Sentence Embeddings using Siamese BERT-Networks,'' in \textit{Proc. EMNLP-IJCNLP}, 2019.
\bibitem{casola2021wits} S. Casola and A. Lavelli, ``WITS: Wikipedia for Italian Text Summarization,'' in \textit{Proc. CLiC-it}, 2021.
\bibitem{beltagy2020longformer} I. Beltagy, M. E. Peters, and A. Cohan, ``Longformer: The Long-Document Transformer,'' \textit{arXiv preprint arXiv:2004.05150}, 2020.
\bibitem{he2022ctrlsum} J. He, W. Kryscinski, B. McCann, N. Rajani, and C. Xiong, ``CTRLsum: Towards Generic Controllable Text Summarization,'' in \textit{Proc. EMNLP}, 2022.
\bibitem{xu2020query} Y. Xu and M. Lapata, ``Coarse-to-Fine Query Focused Multi-Document Summarization,'' in \textit{Proc. EMNLP}, pp. 3632--3645, 2020.
\bibitem{goyal2022news} T. Goyal, J. J. Li, and G. Durrett, ``News Summarization and Evaluation in the Era of GPT-3,'' \textit{arXiv preprint arXiv:2209.12356}, 2022.
\bibitem{dettmers2022llm} T. Dettmers, M. Lewis, Y. Belkada, and L. Zettlemoyer, ``LLM.int8(): 8-bit Matrix Multiplication for Transformers at Scale,'' in \textit{Proc. NeurIPS}, 2022.
\end{thebibliography}

\end{document}
